\hypertarget{operating-events}{%
\section{Operating events}\label{method-corporate-methods-operating-events}}

The operating methods use the selected revenue, inventory and asset provisions of \citet{aasb15dated,aasb102,aasb116dated}. The lease and gross-tax case uses the dated lease, income-tax and cash-flow compilations \citep{aasb16dated,aasb112dated,aasb107dated}. Each rule still needs the entity, period, transaction facts and elections described below.

Every event uses the event-basis, authorization and commit contracts described in Section~\ref{sec:posting-permission}. The event inputs described here are required in addition to identity and date. Each object identifier resolves to its contract, context and admitted source or assumption. A purported cash collection does not establish a sale; a purchase order does not establish receipt.


Revenue, billing and collection are separate transitions. A delivery event names
the contract, performance obligation, quantity, allocated price and control
evidence. Recognition first releases a related contract liability; earned
consideration not yet unconditional creates a contract asset. Both credit
revenue. A missing control finding or unresolved price constraint leaves the
recognition unresolved.

An invoice or right-to-payment event names the contract, bill and unconditional
amount. It debits the receivable and credits an existing contract asset, with
any excess recorded as a contract liability. A collection names that receivable,
settlement and amount and debits cash while crediting the receivable. An excess
payment needs a separate customer-advance method; it cannot be credited to
revenue merely because cash arrived. These transitions retain cumulative
performance, billing and collection independently.

Purchasing starts when the company receives owned goods, not when it places an
order. The event records lot, quantity, unit cost and ownership evidence, debits
inventory and credits the payable. Consigned or unreceived goods cannot pass
that ownership condition. Consumption or sale identifies an available lot
allocation, quantity and admitted cost formula; it debits cost of goods sold
or an eligible production asset and credits inventory. Overselling does not
create stock. Supplier payment separately debits the payable and credits actual
cash. Insufficient cash blocks that settlement and can trigger a separately
modelled default; it never creates a hidden loan.

A tangible-asset event records component cost, residual value, useful life,
depreciation method and the date it becomes available for use. Initial debit to
PPE has a corresponding cash or payable credit; depreciation later debits
expense or an eligible asset and credits accumulated depreciation. Employee
service records earned pay, deductions and employer obligations for the service
period, separately from settlement. Unknown labour or tax law cannot become a
zero payable.

Closing derives revenue, expenses, OCI and owner movements from those events
using admitted statement mappings. It checks cash, inventory, financing and tax
roll-forwards rather than introducing a balancing entry. These event fields
form the method's required inputs; each identifier resolves to the same
contract, context and evidence used by its authorization.


AASB 15 paragraphs 9, 22, 31, 35--43, 47--59, 73--86 and 105--108 provide the selected revenue conditions; the retained compilation dated 31 December 2022 applies to annual periods beginning on or after 1 January 2023 and before 1 July 2026, covering the stated for-profit FY 2024 example. Variable consideration is estimated under a selected method and admitted only subject to the significant-reversal constraint, not recognized solely because its unconstrained expectation is positive. Principal versus agent, returns, warranties, contract modifications and distinct bundles require their own scoped decision records. If their facts are unresolved, retain alternative eligible methods; do not borrow the simple delivery rule by label.

The following schedules make the distinction between performance and billing explicit. For one performance obligation with price 100, cumulative revenue 60, unconditional bills 50 and collection 30, AR=20, CA=10, CL=0 and cash=30. Revenue 60 is represented by those three assets. Billing before performance changes intermediate balances but not the final result. For bills 80 instead, AR=50, CA=0, CL=20; net assets remain 60. Cumulative performed-service cost 40 then 70 posts 40 then 30, not 110. A downward revision follows the correction method and keeps the old result tied to its original view.

For this simple contract, let $R$ be cumulative recognised revenue, $B$
unconditional billings and $C$ collections, with $0\leq C\leq B$. The unpaid
bills determine accounts receivable (AR); performance in excess of billing is
a contract asset (CA), while billing in excess of performance is a contract
liability (CL). After each event,
\[
 \mathrm{AR}=B-C,\quad \mathrm{CA}=\max(R-B,0),\quad
 \mathrm{CL}=\max(B-R,0),\quad \text{cash received}=C.
\]
This recurrence applies per contract, with no impairment, variable consideration or unbilled prepayment; it is not permission to offset unrelated contracts. For billing 80 before delivery, debit AR 80/credit CL 80. Recognizing service 60 later debits CL 60 and credits revenue 60. Collecting 30 debits cash 30/credits AR 30, leaving AR 50 and CL 20, without a second receivable or an artificial contract asset. If service comes first, recognize CA 60/revenue 60, then bill AR 80 against CA 60 and CL 20. Both orders end at the same state. Cash received before an unconditional bill needs an explicit advance event that debits cash and credits CL and records the advance applied to later billing; it cannot pass through collect(AR).

\subsection{A payment before the bill must be applied only once}
\label{sec:customer-advance}

The preceding recurrence deliberately excludes an advance before billing.
Suppose a customer pays 30 towards a service priced at 100. No service has yet
been performed and no payment was already due. The company now has cash 30 and
an obligation to provide the service, so it debits cash and credits a contract
liability. Crediting revenue would balance the entry but recognise performance
that has not occurred. Paragraphs 105--108 of the retained AASB 15 distinguish
this payment from an unconditional receivable \citep{aasb15dated}.

Now let the company earn 60 before issuing a bill. Recognition releases the
liability of 30 and creates a contract asset of 30, with a credit to revenue
of 60. The later invoice shows a gross amount of 80 and applies the customer's
advance of 30, leaving a new unconditional right to 50. Its entry debits the
receivable by 50, credits the contract asset by 30 and credits a contract
liability by 20. Creating a receivable of 80 would count the prepaid 30 again.
When the customer pays another 20, cash is 50, the receivable is 30 and the
contract liability is 20. Their net contribution to assets is the earned 60.

The method therefore keeps the invoice's gross amount separate from the right
still owed. Let $G$ be cumulative gross billings, $A$ cash advances, $K$ the
amount of those advances applied to bills, and $C$ later collections of
receivables. Each application names both its advance and its invoice; the
remaining unapplied credit is $A-K$. As before, $R$ is cumulative earned
revenue. The consideration already paid or made unconditionally due, counted
once, is $T=G+A-K$. For this one fixed-price contract,
\[
 \mathrm{AR}=G-K-C,\qquad
 \mathrm{CA}=\max(R-T,0),\qquad
 \mathrm{CL}=\max(T-R,0),\qquad
 \text{cash received}=A+C.
\]
The minus $K$ prevents the applied advance from appearing both in the cash
payment and in the bill. Since $\mathrm{CA}-\mathrm{CL}=R-T$, the control
\[
 (A+C)+(G-K-C)+\mathrm{CA}-\mathrm{CL}
 =G+A-K+R-T=R
\]
recovers the earned amount. This derivation uses a single contract, with
$0\leq K\leq\min(A,G)$, $0\leq C\leq G-K$ and $0\leq R,T\leq100$.
It provides no authority to offset different contracts or to ignore impairment,
financing, tax, cancellation or refund terms.

Each event must carry its contract, identity, date, currency and admitted
facts. In this restricted extension, advances precede all performance and
billing; later payments settle identified receivables. A revenue increment
first releases an existing contract liability and then creates a contract asset.
For a new gross bill $g$ that applies unused advance credit $k$, only $g-k$
creates a receivable. That amount releases the existing contract asset first;
any remainder increases the contract liability. The invoice consumes its
identified credit even if the advance has already been earned. A duplicate
credit application, an overcollection or missing unconditional-right evidence
cannot pass by adjusting a balancing account.

Billing before performance gives the same final balances in the example but
a different intermediate history. The invoice creates a receivable of 50 and
adds 50 to the existing contract liability of 30; earning 60 then releases 60
of that liability. After the later collection of 20, both paths have the
balances above. Completing the service adds revenue 40, releasing the remaining
liability of 20 and creating a contract asset of 20. A final bill for 20 and
collection of the remaining 50 leave cash 100 and no contract or receivable
balance. The actual sequence is retained even where endpoint amounts agree.

A wholly prepaid invoice exposes another implementation requirement. If its
gross amount and applied advance are both 30, it creates no new monetary right
and needs no accounting lines. It still consumes an identified credit. The
proposed authored method therefore declares a state-only journal proposal:
empty lines, zero account changes and explicit changes to the invoice and
advance allocations. Authorization checks that the net new right is zero;
the same atomic commit and replay rules protect the allocation from reuse.
An empty journal cannot carry a nonzero monetary effect. The native acquisition
schema requires nonempty lines and does not implement this case; an unavailable
authored schema returns an unsupported-operation result without changing state.

These are revenue and cash contributions before service costs and tax. Their
missing measurements remain inputs to the complete company statements. The
next example retains a different obligation: cash already received may have
to be returned even after the goods have been delivered.

\subsection{A return right survives the sale}
\label{sec:customer-returns}

Consider 100 goods sold for 10 each, with a former inventory cost of 6 each.
The buyer obtains control, pays cash and has a contractual right to a refund
for returned non-defective goods. Assume an eligible Australian for-profit
contract in 2024, within the retained AASB 15 compilation's application period.
The contract, control, price-estimation and significant-reversal conditions
are separately satisfied; immediate cash receipt alone would not establish
them. The numbers are synthetic measurements, not estimates of Apple's returns.

If the company initially expects 10 goods to be returned, it expects to retain
900 of the 1,000 received. AASB 15 B20--B25 requires revenue for the expected
entitlement, a refund liability and a separate asset for the right to recover
the goods \citep{aasb15dated}. With no expected recovery costs or deterioration
in this example, the refund liability is 100 and the recovery asset is 60.
The company debits cash 1,000, credits revenue 900 and credits the refund
liability 100. It removes inventory 600 by debiting cost of sales 540 and the
recovery asset 60. Pretax profit is therefore 360. Recognising the full 1,000
as revenue would omit a known refund obligation, although the entry balances.

For a homogeneous cohort of $N$ units, unit price $p$, former unit cost $c$ and
an admitted estimate $q$ of returns, the initial no-recovery-cost case gives
\begin{align*}
 \text{revenue}&=(N-q)p, & \text{refund liability}&=qp,\\
 \text{cost of sales}&=(N-q)c, & \text{recovery asset}&=qc.
\end{align*}
The estimate and its measurement date are retained with the cohort. Paragraphs
53--59 require an appropriate estimation method and the constraint on a
significant reversal; $q$ is not a probability supplied by the accounting
standard. Missing evidence for that measurement leaves the amount unresolved.
After actual returns begin, the state separately records refunded units,
remaining expected returns, remaining liability and recovery rights. A later
outcome cannot be inserted into the original information view.

Suppose the sale occurs in 2024 and new information at year end raises expected
returns to 12. The liability increases by 20, reducing revenue by 20. The
recovery asset increases by 12, reducing cost of sales by 12. Closing revenue
is 880, cost of sales 528 and pretax profit 352. The change in expectations
does not itself establish an error in the original estimate. If the opening
inventory of 600 was financed entirely by opening equity of 600, the resulting
pretax balance sheet has cash 1,000 and recovery asset 72, against refund
liability 120 and equity 952.

In 2025, assume the return period ends with eight accepted returns and no
remaining legal right or unresolved claim. Before final settlement, remeasure
the liability from 120 to 80 and the recovery asset from 72 to 48. The first
change credits revenue by 40; the second debits cost of sales by 24. Refunding
80 then debits the liability and credits cash. Receiving the eight goods
debits inventory by 48 and credits the recovery asset by 48. Settlement does
not recognise the sale or its cost a second time. The 2025 pretax profit
contribution is 16, closing cash is 920 and inventory is 48. With no other
movements, equity is 968: opening 600 plus cumulative profit 368.

The cash and accrual statements can now be checked independently. Operating
cash is 1,000 in 2024 and $-80$ in 2025. In the indirect method, the first
year adds the inventory reduction and liability increase to profit while
subtracting the recovery-asset increase. The second year reverses the relevant
balance movements:
\begin{align*}
 \mathrm{CFO}_{2024}&=352+600-72+120=1{,}000,\\
 \mathrm{CFO}_{2025}&=16-48+72-120=-80.
\end{align*}
There is no investing or financing cash flow, OCI or subsequent owner
transaction in this slice. Tax has not been computed: these pretax
contributions must enter the separate current/deferred-tax method before they
can become complete after-tax company statements.

Recovery costs change the asset without necessarily changing the refund. For
the 12-unit estimate, expected recovery costs of 1 per unit and deterioration
of 0.5 give a recovery asset of $12(6-1-0.5)=54$. Cost of sales becomes 546
and pretax profit 334, while the refund liability remains 120. The cost inputs
are measurements, not automatic zeros; estimating them does not pay cash early.
Paragraph B25 also requires separate presentation of the asset and liability.
Showing only their net 48 in the original no-cost case would conceal both
the gross right to recover goods and the gross obligation to refund cash.

Scope conditions matter as much as the arithmetic. Paragraph B26 treats an
exchange for the same type, quality, condition and price differently from a
return; B27 sends a defective-product replacement to warranty guidance.
The refund method cannot accept those events simply because the goods move
back to the seller. It also needs the contractual return window, accepted-unit
identities, actual refund and recovery measurements, and evidence that any
remaining rights have expired before clearing the balances. An unknown
estimate yields unresolved, absent control makes recognition ineligible, and
an event of a different kind remains outside this method. Inventory costing
and the condition of recovered goods continue to have their own requirements.

\subsection{A warranty can promise another year of service}
\label{sec:service-warranty}

A returned product and a warranty claim need different histories. The return
example unwinds a sale when a customer exercises a refund right. A warranty
may instead promise to keep a product working, or provide a service after the
product has been delivered. The company must identify which promise it made
before deciding how much of a customer's payment it has earned. The following
cases use the same retained AASB 15 compilation and historical Australian
for-profit scope as the return example.

Suppose a customer buys a product and a separately purchasable service warranty
for a combined 108. Comparable standalone selling prices at contract inception
are 100 for the product and 20 for the warranty. Both are distinct promises;
the product works independently of the additional service. The contract has
neither a discount attributable to just one promise nor variable consideration.
Under paragraphs 73--86, the 12 bundle discount is therefore allocated in
proportion to the standalone prices. The product receives
$108\times100/120=90$ and the service receives $108\times20/120=18$.
An invoice's labels would not by themselves establish those standalone prices
or justify a different discount allocation \citep{aasb15dated}.

On 1 July 2024 the company receives 108 and transfers control of the product,
whose inventory cost is 60. First, receipt debits cash and credits the contract
liability by 108. Delivery then debits that liability and credits product
revenue by 90, while debiting cost of sales and crediting inventory by 60.
Profit is 30 and a liability of 18 remains. Recognising all 108 on delivery
would credit revenue for service the company still owes.

The service warranty runs for twelve monthly service periods to 30 June 2025.
The company stands ready uniformly; the customer receives and consumes the
benefits as service is supplied. The admitted measure gives each completed
service month one twelfth of the allocation and faithfully depicts performance.
These facts support over-time recognition under
35(a), 39--45 and B14--B19; the existence of a twelve-month date range alone
would not suffice. At 31 December 2024, half the service has been provided,
so half of 18, or 9, is earned. Debit the contract liability and credit
service revenue by 9. Actual service costs of 3 are expensed and paid.
Opening inventory and equity were both 60, with no other opening balances.
Year-one revenue is now 99, expenses 63 and pretax profit 36. Cash 105 equals
the remaining liability 9 plus equity 96.

The second half-year releases the remaining liability of 9 and incurs another
3 of paid service costs. Its profit is 6 and closing cash and equity are both
102. Over the complete contract, revenue 108 less costs 66 gives profit 42.
The cash statement follows the same history: operating cash is 105 in the
first period and $-3$ in the second. The independent accrual bridges are
$36+60+9=105$ and $6-9=-3$, using the inventory reduction and changes in the
contract liability. There is no OCI or subsequent owner transaction in this
pretax slice; the complete company calculation still needs tax and any other
obligations.

The contract record must therefore retain the two promises, the original
allocation and its price evidence, delivery identity, coverage, progress
measure, cumulative service revenue and actual costs. Each service update
posts the difference between newly supported cumulative revenue and revenue
already recognised. Replaying an accepted delivery cannot release the product
allocation twice. Unknown progress is not zero performance: paragraphs 44--45
require a separate assessment of reasonable measurement and, where applicable,
cost recovery. This percentage method cannot invent those findings.

The classification also has a substantive boundary. A separately purchasable
warranty is a distinct service under B29. Without that option, B30 sends an
assurance-only warranty to AASB 137 unless it provides additional service.
Law, duration and promised tasks help distinguish the cases under B31; there
is no universal number of months that makes a warranty a service. If assurance
and additional service cannot reasonably be accounted for separately, B32
treats the two warranties together as one performance obligation. Harm or
infringement compensation under B33 is not such a service. The exact Australian
AASB 137 measurement text has not been inspected for this extension, so an
assurance obligation retains a required measurement method and facts. It is
not assigned a zero provision or converted into the service example.

\subsection{Changed promises can affect future units or past progress}
\label{sec:customer-modification}

A contract can change after some performance has already been recognised.
The accounting consequence depends on what remains to be supplied. Consider
an original prepaid order for ten distinct units at 10 each. Four have been
delivered, earning 40. The parties approve an enforceable amendment adding
six units for 48. Comparable standalone prices remain 10 each, and no
appropriate contract-specific adjustment makes the 48 addition reflect
those prices. The addition therefore fails the separate-contract condition
in paragraph 20. Because the remaining units are distinct from those already
delivered, paragraph 21(a) reallocates the consideration for the remaining
performance \citep{aasb15dated}.

The undelivered original units carry 60 of consideration. Adding the new 48
gives 108 for twelve remaining units, whose standalone prices are equal.
Thus their new unit allocation is
\[
 \frac{\text{unrecognised original consideration}
       +\text{additional consideration}}
      {\text{remaining equal-price units}}
 =\frac{60+48}{6+6}=9.
\]
The first four units keep their revenue of 40. Spreading the entire revised
price over all sixteen units would instead give $148/16=9.25$ and rewrite
their revenue as 37. That calculation is wrong for the stated facts: this
modification reallocates remaining performance, not completed deliveries.

Inventory and cash make the consequences concrete. Start with sixteen units
costing 6 each, financed by opening equity of 96; six units were initially
uncommitted. The first prepayment debits cash and credits a contract liability
by 100. The first four deliveries release 40 of revenue and remove inventory
cost 24. Immediately before amendment, cash 100 plus inventory 72 equals the
liability 60 plus equity 112. Receipt of the additional 48 increases cash
and the liability. Approving the changed order itself earns no new revenue
and creates no new stock.

Deliver six more units by the first year end. They release liability and
revenue of 54, with inventory cost 36. Total year-one revenue is 94, cost
of sales 60 and profit 34. Cash 148 and inventory 36 equal liability 54
and equity 130. Delivering the final six in the next year earns another
54 at cost 36, leaving cash and equity 148 with no inventory or contract
liability. That period's profit is 18 but no cash arrives: the customer
already paid. The operating cash bridges are $34+60+54=148$ and
$18+36-54=0$. Lifetime revenue 148 less inventory cost 96 gives profit 52.

If the additional six units instead cost 60 and meet paragraph 20's
standalone-price condition, the addition is a separate contract. The first
four still earn 40; the six remaining original units earn 60 and the new
six earn 60. Both contracts keep their own allocation records. The difference
between these cases is the admitted price and scope assessment, not permission
to choose whichever result produces the preferred margin.

An integrated promise can have the opposite effect on past cumulative
revenue. Suppose a customer controls work in progress as the company creates
it, satisfying paragraph 35(b). The company has one integrated performance
obligation, and its remaining work is not distinct from work already supplied.
Original price 100 and supported output progress of 40\% yielded revenue 40.
An approved enforceable amendment changes that same promise to price 130;
the revised faithful output measure is 30\% of the expanded work. Under
21(b), both the price and progress change the cumulative amount earned:
\[
 \text{revised cumulative revenue}=130\times0.30=39,
 \qquad \text{catch-up}=39-40=-1.
\]
The new work makes the completed work a smaller share of the promise.
Revenue decreases by 1 at amendment even though the total contract price
increases. The old 40 remains attached to its original information view;
the modification does not establish that the earlier estimate was an error.

For a complete prepaid example, start with cash and equity 50. Receive 100
and pay service costs 30 before amendment. Recognition of 40 leaves cash
120, liability 60 and equity 60. At amendment, receive the additional 30
and debit revenue 1 while crediting the liability 1. Cash is now 150,
the liability is 91 and equity is 59: cumulative revenue 39 less costs 30
has earned profit 9. Operating cash since opening is 100, also $9+91$.
Completing the promise next period releases revenue 91 and incurs further
paid service costs of 60. Period profit is 31, operating cash is
$31-91=-60$, and closing cash and equity are 90. Lifetime profit is
$130-90=40$. The output progress percentages were independently supported;
they were not inferred from these cash costs.

Both modifications need the approved rights, effective date, previous and
revised promise, consideration, performance measurements and source versions.
The distinct-goods method freezes the fulfilled deliveries and replaces the
remaining allocation; the integrated method proposes only the difference
between revised cumulative revenue and the amount already recognised.
Economic state and journals change atomically. Different billing schedules
can require a contract-asset adjustment instead of the prepaid liability
entries used here. A request that has not been approved leaves the existing
contract operative under paragraph 18. Unknown enforceability stays
unresolved; approved scope with undetermined price requires paragraph 19's
variable-consideration estimate and constraint, not a price of zero.

Paragraph 21(c) also admits mixed remaining obligations, requiring an explicit
allocation consistent with the two principles above. The examples do not
force such a contract into one category. Price-estimate changes without a
contract modification follow 87--89. A later price change after modification
needs paragraph 90: variable consideration promised before a 21(a)
modification can still be allocated to the earlier obligations under 90(a).
It cannot automatically be spread over the twelve remaining units merely
because their current allocation is 9 each.

\subsection{The same cash can accompany different revenue}
\label{sec:principal-agent-cash}

Suppose an entity receives 120 from a customer and sends 100 to a supplier.
It keeps 20. Those cash movements cannot tell us whether it sold the supplier's
good or arranged a sale by that supplier. Under B34--B38, the decision starts
with the specified good or service and whether the entity controls it before
transfer to the customer. Responsibility for delivery, inventory risk and
price discretion inform that decision; none is a substitute for control, and
momentary legal title does not by itself establish it \citep{aasb15dated}.

In the principal case, the entity controls the identified good before
customer transfer, can direct its use and obtain its benefits, and is
responsible for delivering it. Acquisition on credit debits inventory 100
and credits the supplier payable 100. Customer transfer and cash receipt
debit cash 120 and credit revenue 120; the inventory leaves through cost
of sales 100. Settling the supplier debits the payable and credits cash 100.
From zero opening balances, the result is cash and equity 20, revenue 120
and profit 20.

In the agent case, the supplier controls the good until transfer; the entity
only arranges the sale and earns a fixed commission of 20 when that service
is completed. Assume the entity controls the collected cash and owes the
supplier 100 under its remittance agreement. Funds held outside its balance
sheet in a fiduciary capacity are outside this example. Receipt therefore
debits cash 120, credits the supplier payable 100 and credits commission
revenue 20. Settlement clears the payable against cash 100. The agent has
no inventory or supplier cost of sales. Its final cash, equity and profit
are again 20, but revenue is 20.

Even the intermediate balance sheets agree: before remittance, both have
cash 120, a payable 100 and equity 20. The different control facts explain
the different revenue. The method must retain those facts and classify each
specified promise separately, since the same entity can act as principal
for one promise and agent for another. An unknown control finding leaves
classification unresolved; the receipt cannot decide it. Both cash legs
remain in the ledger. Whether the cash-flow statement presents receipts and
remittances gross or net requires its own AASB 107 conditions, which do not
follow merely from the revenue classification. These examples establish
the cash change and pretax profit contribution while preserving that separate
presentation question.


\subsection{Inventory and assets retain their cost histories}
\label{sec:corporate-cost-histories}

Inventory uses identified cost for non-interchangeable items or the admitted cost formula for interchangeable items. AASB 102.9, 23, 25 and 28--34 support lower of cost/net realizable value, relevant cost formulas and write-down/reversal conditions. A method never chooses LIFO merely to fit a reported balance. In a second lot example, purchases are 10 units at 2 and 10 at 3, followed by sale of 12 units. FIFO gives cost 26 and stock 8 units/cost 24. Weighted average gives cost 30 and stock 20. Under the same given selling price 48, margin differs 22 versus 18. Both alternatives require their own permitted policy, consistent application and evidence; matching total cash does not remove the difference. NRV 16 for the remaining FIFO lot records expense 8, inventory allowance 8; recovery to NRV 22 reverses 6, leaving allowance 2. Recovery to 30 is capped at original cost 24. US LIFO remains a separately sourced question.

The retained AASB 116.7/15, 29--30 and 48--62 support cost and depreciation. Asset 100, residual 10, three-year even use gives 30 per full year. A year-two estimate change after first-year depreciation to residual 10 and three remaining years gives (70\ensuremath{-}10)/3=20 prospectively; an error in original cost is a correction, not an estimate change. Component replacements derecognize the replaced component's remaining amount before recognizing its replacement; no perpetual double asset. Borrowing-cost capitalization, impairment, provisions, employee benefits and share-based payment are specified as additional scoped methods below, not silently assigned zero in real-company cases.

\section{Apple: asset balances do not identify asset movements}
\label{sec:apple-operating-assets}

The preceding methods need receipts, lots, useful lives and settlement dates.
Apple's disclosures supply some of the balances those events must explain,
but they do not supply the events themselves. Its FIFO and straight-line
policies narrow the permitted cost calculations without identifying the actual
inventory lots or asset cohorts \citep[Note 1, printed p.~34]{apple2024}.
Comparing the notes with the cash-flow statement shows why filling those
missing inputs by subtracting balances can produce a wrong history.

\subsection{An inventory adjustment is not purchases or payments}

Between 28 September and 28 December 2024, components fall from 3,627 to 2,792
million dollars while finished goods rise from 3,659 to 4,119 million. Total
inventory falls from 7,286 to 6,911: the component changes, $-835$ and $+460$,
sum to $-375$ million \citep[Note 5, printed p.~9]{apple2025q1}. Those two
classes do not reveal quantities, unit costs or the transfers and sales that
produced the balances.

The quarter's cash-flow statement reports $+215$ million for inventories within
its reconciliation of net income to operating cash \citep[printed p.~5]{apple2025q1}.
That is an indirect reconciliation adjustment, not total cash paid to suppliers.
Even replacing it with minus the stock change would give $+375$, a different
quantity. The difference $215-375=-160$ million remains unexplained by these
selected disclosures. In fiscal 2024, inventory rose from 6,331 to 7,286, while
the cash-flow adjustment was $-1{,}046$; relative to the $-955$ stock-change
shortcut the difference is $-91$ million \citep[printed pp.~31 and 33]{apple2024}.
Neither residual is an identified currency effect, write-down or other journal
entry. Production, consumption, ownership, valuation and payment histories must
be established before they can explain the differences.

\subsection{Testing a proposed fixed-asset history}

The fixed-asset notes distinguish three measurements. The table shows gross
property, plant and equipment (PP\&E), its accumulated depreciation and the
remaining carrying amount, all in millions of dollars
\citep[Note 5, printed p.~39]{apple2024}
\citep[Note 5, printed p.~9]{apple2025q1}.

\begin{center}
\begin{tabularx}{\linewidth}{@{}Xrrr@{}}
\toprule
Amount & 30 Sep 2023 & 28 Sep 2024 & 28 Dec 2024\\\midrule
Gross PP\&E & 114,599 & 119,128 & 120,615\\
Accumulated depreciation & 70,884 & 73,448 & 74,546\\
Net PP\&E & 43,715 & 45,680 & 46,069\\\bottomrule
\end{tabularx}
\end{center}

Accumulated depreciation is shown here as a positive deduction: gross less
that deduction equals net. During fiscal 2024, gross cost increases 4,529,
accumulated depreciation increases 2,564 and net PP\&E increases 1,965.
Yet the note reports PP\&E depreciation expense of 8.2 billion, and the
cash-flow statement reports capital payments of 9,447 million. Its broader
depreciation-and-amortisation adjustment is 11,445 million
\citep[printed pp.~33 and 39]{apple2024}. These amounts differ in meaning;
the 8.2 billion also has coarser displayed precision. The change in accumulated
depreciation cannot be substituted for the year's depreciation charge.

Consider a deliberately limited asset book to test a tempting reconstruction.
Let $G_0,G_1$ be opening and closing gross cost and $B_0,B_1$ accumulated
depreciation, recorded as positive amounts. Let $A$ be gross additions, $R$
the original cost of retired assets, $d$ the depreciation charge and $S$ the
accumulated depreciation removed with those assets. Assume there are no other
movements and all depreciation is expensed. For this illustration, use the
displayed amounts as exact, including $d=8{,}200$ million. The two balances
must then satisfy
\begin{equation}
 G_1=G_0+A-R,\qquad B_1=B_0+d-S,\qquad 0\leq S\leq R.
 \label{eq:apple-asset-movement}
\end{equation}
The inequality requires a nonnegative carrying amount for the retired assets
in this cost model. If additions are equated with the 9,447 of cash payments,
the gross balance requires $R=9{,}447-4{,}529=4{,}918$. The depreciation
balance requires $S=8{,}200-2{,}564=5{,}636$. The proposed book removes more
accumulated depreciation than original cost, by 718 million. It is infeasible
under these combined assumptions. That finding rejects the proposed history;
it does not establish an error in Apple's filing or identify which assumption
must change.

The result is not sensitive to small rounding differences under a specified
rounding model. If the million-dollar balances and payment are rounded to the
nearest million, and 8.2 billion to the nearest 0.1 billion, the excess
$d-9{,}447+(45{,}680-43{,}715)$ lies in the closed outer envelope
$[666.5,769.5]$ million. Its lower endpoint is
$8{,}150-9{,}447.5+45{,}679.5-43{,}715.5=666.5$.
This is a conditional precision check, not a verified issuer rounding rule;
the exact tie convention remains unknown.

Allowing additions to be partly unpaid provides two simple repairs. Retain
the same selected opening and closing balances, 8,200 depreciation and 9,447
cash payments. Assume retirement at the start of the year with zero proceeds,
depreciation of surviving assets, and new acquisitions at year end. There are
no opening capital payables, taxes or other asset movements. The figures below
are constructed histories, not estimates of actual Apple transactions.

\begin{center}
\begin{tabularx}{\linewidth}{@{}Xrr@{}}
\toprule
Annual movement (USD million) & History 1 & History 2\\\midrule
Gross additions & 10,529 & 11,529\\
Gross cost retired & 6,000 & 7,000\\
Accumulated depreciation removed & 5,636 & 5,636\\
Retired carrying amount, expensed & 364 & 1,364\\
Cash paid for additions & 9,447 & 9,447\\
Additions remaining payable & 1,082 & 2,082\\\bottomrule
\end{tabularx}
\end{center}

Both histories satisfy equation~\eqref{eq:apple-asset-movement}; their gross
additions can also be allocated to match the three reported asset classes.
For History 1, acquisition debits PP\&E cost 10,529 and credits cash 9,447 and
a capital payable 1,082. Depreciation debits expense and credits accumulated
depreciation by 8,200. Retirement debits accumulated depreciation 5,636 and
loss 364, while crediting gross PP\&E cost 6,000. Net PP\&E rises 1,965 and
cash falls 9,447, so total assets fall 7,482. Liabilities rise 1,082 and equity
falls 8,564 through the two expenses: $1{,}082-8{,}564=-7{,}482$.
History 2 has the same asset change but liabilities rise 2,082 and equity
falls 9,564. These are coherent subsystem consequences, not a fit to all of
Apple's statements or an exhaustive family of real asset histories. Neither
history receives a probability or preference from this arithmetic.

The following quarter supplies a new snapshot, not the missing annual movements.
Its gross, accumulated and net changes are 1,487, 1,098 and 389 million,
whereas capital payments are 2,940 and broad depreciation/amortisation is
3,080 \citep[printed pp.~5 and 9]{apple2025q1}. The inspected interim note does
not separately identify PP\&E depreciation expense. The lease note adds a
related distinction: finance-lease right-of-use assets are included in PP\&E,
while operating-lease right-of-use assets are in other noncurrent assets.
The 1.0 billion of noncash right-of-use acquisitions covers the disclosed lease
population \citep[Note 8, printed p.~42]{apple2024}; adding it all to PP\&E
or to cash capital expenditure would erase that distinction.

The method therefore preserves reported balances, cash payments, expense and
unresolved movement components separately. A future cost or funding query
still needs lot and asset ages, useful lives, residual values, impairment and
settlement facts. The real filing has supplied constraints and ruled out one
simple history, while leaving several different explanations open.

\subsection{Local alternatives need different company histories}

The two asset books expose a constraint on combining explanations. Both fit
the selected PP\&E movements and capital payments. Can they therefore be
inserted interchangeably into one otherwise unchanged history of Apple?
The journals above answer no. Relative to History 1, History 2 retires an
additional 1,000 of gross cost without removing additional accumulated
depreciation, so it recognises an additional 1,000 loss. It acquires an
additional 1,000 of gross cost on credit. Those changes cancel in net PP\&E
and leave cash unchanged, but they add 1,000 to liabilities and subtract
1,000 from profit and equity. Amounts remain millions of dollars, under the
same conditional asset-book assumptions, including exact illustrative
depreciation of 8,200 and no tax effect within that subsystem.

Let $P_i$ denote the vector of statement effects generated by asset History
$i$, and let $R_i$ denote the effects of the rest of the company's history.
For this comparison, order each vector as asset change, liability change,
profit, equity change and cash change. Both full histories must explain the
same reported vector $Y$. If an admissible remainder exists for each branch,
their necessary relation is
\[
P_1+R_1=Y=P_2+R_2,
\qquad
R_2-R_1=-(P_2-P_1)=(0,-1{,}000,+1{,}000,+1{,}000,0).
\]
Thus the rest of History 2 must contribute 1,000 less liability growth and
1,000 more profit and equity growth, with no additional asset or cash change.
Holding the remainder fixed is wrong relative to the common-statement target.
The balance-sheet identity alone fails to detect the error: adding 1,000 to
liabilities and subtracting 1,000 from equity still leaves their sum unchanged.

The reported annual totals make the required differences concrete. Apple's
asset growth is 12,397, liability growth is 17,593, profit is 93,736 and equity
change is $-5{,}196$ \citep[printed pp.~29--32]{apple2024}. Since History 1
contributes liability growth of 1,082, profit of $-8{,}564$ and equity change
of $-8{,}564$, its hypothetical remainder must contribute 16,511, 102,300
and 3,368 respectively. History 2 instead requires 15,511, 103,300 and
4,368. In both cases the remaining asset change must be 19,879. These are
necessary subtraction results. They do not supply the transactions, expense
classifications, taxes or liability assignments that would make a remainder
admissible.

The distinction matters for cash as well. In the isolated asset book,
depreciation and the retirement loss are noncash expenses. Adding both back
to its negative profit gives zero operating cash; acquisitions use 9,447 of
investing cash. Switching branches therefore adds 1,000 to the noncash loss
and 1,000 to the corresponding addback, leaving subsystem cash unchanged.
If that loss were mapped into the annual cash-flow statement's broad Other
adjustment of $-2{,}266$, the remaining contributions to that row would have
to be $-2{,}630$ in History 1 and $-3{,}630$ in History 2
\citep[printed p.~33]{apple2024}. That mapping is itself conditional: the
source does not identify either retirement loss or the composition needed
to offset it. A matching operating-cash total cannot substitute for the
reported adjustment rows.

Nor may a remainder erase the cash-population difference from
Section~\ref{sec:apple-statements}. After the common asset payments, the
required remaining balance-sheet cash change is 9,425, whereas the required
remaining cash-flow change is 8,653. Their 772 difference is inherited from
the disclosed opening populations. Calling it another cash receipt would
manufacture a transaction rather than explain the record.

A proposed completion must now identify the liability class containing the
capital payable, the expense class containing the loss, the tax consequences
and the noncash cash-flow mapping, and satisfy the detailed statement rows
under the same assignments. It must also retain the asset cohorts and payment
terms needed at the next date. The December snapshot does not establish that
either hypothetical annual payable remained outstanding or identify when it
was settled. Failure of the unchanged-remainder construction therefore calls
for a coupled historical model. It rejects that construction without deciding
whether either local asset book has an admissible full-company completion.


\subsection{An offset must explain the individual rows}
\label{sec:apple-coupled-completion}

The required offset can be tested without pretending that its transaction has
been observed. Suppose a separate economic obligation would justify an
expense and an unpaid liability of $q$ in History 1 and $q-1{,}000$ in
History 2, with $q\geq1{,}000$. These are alternative amounts in two histories;
they are not a release booked within either history. For a numerical example,
take $q=2{,}000$: each branch debits expense and credits its own liability,
for 2,000 or 1,000 respectively. The difference supplies the 1,000 lower
liability and 1,000 higher profit required of History 2's remainder.
The economic obligation is an additional premise. Subtracting reported
totals has not established that it exists.

Even conditional on that premise, a matched total is insufficient. Assume
both the retirement loss and the new expense enter product cost, and both
the capital payable and the new liability enter the same non-tax current
liability category. History 2's extra loss then offsets its smaller expense
within that expense row, and its larger capital payable offsets its smaller
expense liability within that liability row. Other operations, distributions,
share counts and OCI remain fixed in this comparison. Equal tax expense and
tax balances are a further condition: equal pretax profit does not imply
equal deductions, tax bases or deferred tax.

The indirect cash-flow statement exposes another restriction. If the
expense liability is treated as an ordinary operating accrual, its 1,000
smaller increase reduces the working-liability adjustment by 1,000. The
extra retirement loss increases the noncash addback by 1,000. Total operating
cash flow is unchanged, but the separately reported Other adjustment and
Other current and non-current liabilities adjustment each change. That
construction fails to reproduce Apple's detailed cash-flow statement,
even though it balances and matches total operating cash
\citep[printed p.~33]{apple2024}.

A different conditional reporting map identifies the expense liability's
recognition as a separate noncash adjustment. Its amount must then be
excluded from the working-liability adjustment to avoid adding it back
twice. Let $\alpha$ be the fraction of this hypothetical recognition
allocated to the separate noncash adjustment; the remaining $1-\alpha$
enters the operating-liability adjustment. This parameter describes an
assumed allocation, not a probability or a freely selectable accounting
policy. With the retirement loss also assigned to Other, the differences
between History 2 and History 1 are
\[
\Delta\mathrm{Other}=1{,}000-\alpha\,1{,}000,
\qquad
\Delta\mathrm{WorkingLiabilities}=-(1-\alpha)\,1{,}000,
\qquad
\Delta\mathrm{CFO}=0.
\]
For $\alpha=0$, the two rows differ by $+1{,}000$ and $-1{,}000$.
For $\alpha=\tfrac12$, they still differ by $+500$ and $-500$.
Both row differences vanish only at $\alpha=1$ within this construction.
Including the full liability movement as well as the separate addback
would instead understate the History 2 minus History 1 operating-cash
difference by 1,000.

Under the common expense and liability mappings, equal taxes and the
disjoint noncash mapping with $\alpha=1$, all five annual statements are
unchanged by the combined transformation. Fixed shares also preserve EPS;
fixed OCI and distributions carry the profit equality into comprehensive
income and the equity statement. This is a conditional invariance result.
It assumes a feasible full-company History 1 exists and that the transformed
events remain admissible. It neither constructs that base nor proves that
either complete history is compatible with all notes. Assigning the two
expenses to different rows, the liabilities to different categories, or
different tax consequences immediately breaks the relevant equality.

The notes make these remaining conditions specific. Note 5 separates gross
PP\&E by asset class: the extra acquisition and retirement must cancel
within each affected class to leave that table unchanged. Note 6 separates
26,601 of current income taxes payable from 51,703 of other current
liabilities, and 9,254 of non-current income taxes payable from 36,634 of
other non-current liabilities. A total-liability offset cannot choose its
own location inside those amounts. Note 8 identifies finance-lease
right-of-use assets of 1,069 within PP\&E and distinguishes lease liabilities
and noncash lease acquisitions. Those disclosures do not identify the
hypothetical capital purchases or expense liability. Nor do Note 12's
purchase commitments and contingency assessment establish the assumed
obligation or its amount
\citep[printed pp.~39, 42, 45--46]{apple2024}.

The distinction therefore survives the repair: an algebraic offset is
available; a row-preserving transformation is available under stated
economic, tax and reporting assumptions; an actual Apple completion is
unestablished. To admit one, the historical model needs the underlying
obligation, gross asset classes, expense and liability assignments, tax
bases and noncash treatment, together with a feasible rest of the company.
The subsequent quarter additionally needs settlement dates and surviving
asset cohorts. Identical annual statements do not make those future
responses identical.


\subsection{Equal liabilities can require different cash paths}
\label{sec:apple-settlement-paths}

The conditional repair leaves different obligations beneath the same annual
liability total. In the $q=2{,}000$ illustration, History 1 owes 1,082 for
capital purchases and 2,000 for the separate expense; History 2 owes 2,082
and 1,000 respectively. Each owes 3,082. Consider what paying those balances
would do, conditional on the economic and reporting assumptions needed for
the preceding construction. Assume payment at carrying value, sufficient
unrestricted cash, and no additional interest, foreign exchange,
remeasurement or tax effect. Capital settlement is classified as investing
cash and expense settlement as operating cash in this example. Neither
the obligation nor its settlement terms have been identified in Apple's
disclosures.

If both obligations are paid at the first subsequent date, each payment
debits its liability and credits cash. No new expense arises: the expense
or acquisition was recognised in the prior period. Both histories reduce
cash and liabilities by 3,082, leave assets other than cash unchanged, and
leave profit, OCI and equity unchanged. The indirect operating-cash
calculation includes the expense liability's decrease once; it excludes
capital settlement from the operating-liability adjustment. The
cash-flow statement nevertheless distinguishes the histories. Writing
each ordered pair as operating cash followed by investing cash, the
joint alternatives are
\[
(\mathrm{CFO},\mathrm{CFI})_{H_1}=(-2{,}000,-1{,}082),
\qquad
(\mathrm{CFO},\mathrm{CFI})_{H_2}=(-1{,}000,-2{,}082).
\]
Both pairs sum to $-3{,}082$. Selecting the smaller outflow from each
category would instead produce $(-1{,}000,-1{,}082)$, whose total is only
$-2{,}082$. No branch produces that pair. Selecting both larger outflows
would invent an extra 1,000 payment. The category intervals can describe
marginal possibilities, but their Cartesian product contains paths
excluded by the shared history.

Timing makes the distinction relevant even to total cash. Suppose capital
obligations are paid at the first date and expense obligations at the
second. History 1 then pays 1,082 followed by 2,000; History 2 pays 2,082
followed by 1,000. Relative to History 1, History 2 has 1,000 less cash
after the first date and the same cumulative cash after the second.
Paying the expense obligations first reverses the first-date difference.
An aggregate liability of 3,082 suffices for the narrow query of total
undiscounted cash after both payments, under these fixed assumptions.
It does not suffice for cash by date or by statement category. Retaining
liability kind, amount, settlement terms and joint branch identity is
therefore necessary for those queries.

The actual following quarter constrains a proposed continuation without
identifying it. Apple reports PP\&E payments of 2,940, operating cash of
29,935, investing cash of 9,792 and total cash growth of 356 in the quarter
ended 28 December 2024 \citep[printed p.~5]{apple2025q1}. If both
hypothetical obligations settled in that quarter, the remaining PP\&E
payments would have to be 1,858 under History 1 or 858 under History 2.
The rest of the company would have to contribute operating cash of
31,935 or 30,935, and investing cash of 10,874 or 11,874, respectively.
After unchanged financing cash of $-39{,}371$, either required remainder
has cash growth of 3,438; subtracting the hypothetical 3,082 settlement
recovers the reported 356. These are necessary conditions on unconstructed
remainders. Positive remaining payments and a matching cash identity do
not establish the underlying transactions or prove that either
continuation fits the complete quarterly record.

This calculation extends the conditional history into explicit journals
and statement effects while keeping actual-history admission unresolved.
If cash is insufficient, the model must retain the resulting funding or
default question; it cannot invent a loan to preserve the path. A later
depreciation or production query requires different information again:
asset cohorts, remaining lives, residual values, capacity and tax bases.
The settlement example supplies none of those quantities. A change in the
query must therefore reopen the state-sufficiency check, even when the
annual statements and the final cash total coincide.



\hypertarget{lease-and-tax-integration}{%
\section{Lease and tax integration}\label{method-corporate-methods-lease-and-tax-integration}}

A lease connects asset use, a financing obligation and income tax. The following schedule extends the operating-company example by separating these three measurements. The retained lease and tax dossiers supply the dated provisions, modification case and assumptions. The official lease compilation is dated December 2023 despite the older filename. The right-of-use (ROU) asset represents the admitted use of the leased asset. Its cost, the lease liability, accumulated depreciation, remaining payments, rate version and option judgments are separate state fields. Tax retains carrying amounts, legal tax bases, reversal dates, enacted/substantively-enacted rate evidence, recovery support, and original profit/OCI/equity allocation. The deferred tax asset (DTA) and deferred tax liability (DTL) are retained gross before any permitted offset presentation. A liability schedule is not a tax-base schedule.

For the two annual 121 payments at 10\%, initial liability and ROU are 210. Year-one interest is 21, liability after payment 110, depreciation 105 and ROU 105. Year two interest is 11, and both balances finish at zero. A same-asset modification just after year one creates two future payments 144 at revised 20\%, PV 220, so liability increases 110 and ROU increases 110 to 215. A payment index reset has different rate rules; it cannot reuse that modification blindly.

To demonstrate connection to company statements, assume a synthetic entity with opening cash/equity 500, no other opening balances, and annual cash service revenue 400 less cash operating costs 100, all recognized in their years. There are no dividends, other assets, debt, losses or permanent differences. Add the base lease above. Assumed tax law deducts cash lease payments against taxable income, attributes those deductions to the lease liability, gives zero tax bases to ROU/liability, and taxes profits at enacted 25\%. Sufficient usable profits and all recognition conditions for DTA are assumed. Tax is paid the following year. No current-period tax payment is invented from tax expense.

At commencement debit DTA 52.5 and credit DTL 52.5 in addition to the lease recognition. At year-one close DTA is 27.5 and DTL 26.25. The DTA reduction 25 is deferred expense; DTL reduction 26.25 is deferred benefit, net benefit 1.25. Taxable profit is 300\ensuremath{-}121=179; current tax/payable 44.75. Pretax book profit is 300\ensuremath{-}105\ensuremath{-}21=174, total tax 43.5 and profit 130.5. At year two, deferred expense 1.25 combines with current tax 44.75 to give total tax 46 on pretax profit 184.

\begin{center}
\begin{tabular}{lrr}
\toprule Closing balance & 2024 & 2025\\\midrule
Cash & 679 & 813.25\\
ROU asset & 105 & 0\\
Deferred tax asset & 27.5 & 0\\
Lease liability & 110 & 0\\
Deferred tax liability & 26.25 & 0\\
Current tax payable & 44.75 & 44.75\\
Equity & 630.5 & 768.5\\\bottomrule
\end{tabular}
\end{center}


2024 assets 811.5 equal liabilities 181 plus equity 630.5; 2025 assets 813.25 equal liabilities 44.75 plus equity 768.5. For this case, elect financing classification for lease interest under the retained AASB 107 regime. The operating and financing cash totals below depend on that presentation election. Direct operating cash is 300 then 255.25 after the first tax payment; financing lease cash is \ensuremath{-}121 each year; no investing cash arises. Indirect operating cash in year one is profit 130.5 plus depreciation 105 plus interest 21 minus deferred benefit 1.25 plus unpaid current tax 44.75 = 300. Year two gives profit 138 +105 +11 +1.25 =255.25. The second tax payable settles next year. Zero OCI and no owner transactions mean the equity change equals cumulative profit. This case tests gross tax, cash timing and financial-statement coherence, not a real entity's tax law.

\section{Apple: recognising tax and settling the obligation}
\label{sec:apple-tax-settlement}

The lease example separates a tax charge from the payment that follows it.
Apple's fiscal 2024 disclosures show why that separation matters when several
tax jurisdictions, deferred balances and a large disputed obligation enter
the same accounts. Its annual income-tax expense was 29,749 million dollars;
net cash paid for income taxes was 26,102 million
\citep[printed pp.~29 and 33]{apple2024}. The 3,647 difference is not, by itself,
an increase in tax payable. Deferred tax and other balance movements must first
be distinguished from current tax and settlement.

\subsection{The expense and balance reconciliations answer different questions}

Note 7 divides the 2024 expense between current and deferred amounts for federal,
state and foreign taxes. In millions of dollars, the reported components are
as follows \citep[Note 7, printed p.~40]{apple2024}. Parentheses in the source
represent negative expense, or a tax benefit.

\begin{center}
\begin{tabularx}{\linewidth}{@{}Xrrr@{}}
\toprule
Tax category & Current & Deferred & Total\\\midrule
Federal & 5,571 & $-3{,}080$ & 2,491\\
State & 1,726 & $-298$ & 1,428\\
Foreign & 25,483 & 347 & 25,830\\
Sum & 32,780 & $-3{,}031$ & 29,749\\\bottomrule
\end{tabularx}
\end{center}

The current charge is therefore 32,780 and the deferred benefit 3,031.
The tax payable has its own scope. Note 6 reports current income taxes payable
of 26,601 and noncurrent income taxes payable of 9,254 at 28 September 2024,
against 8,819 and 15,457 a year earlier. Combining both maturity classes gives
35,855 against 24,276, an increase of 11,579
\citep[Note 6, printed p.~39]{apple2024}. Current expense less net cash taxes
is only $32{,}780-26{,}102=6{,}678$. Reclassifying an amount between current
and noncurrent cannot explain the difference in their combined total.

To expose the missing movements, let $T$ denote that combined tax payable and
$N$ the net deferred tax asset reported in Note 7. Let $C$ be current tax
expense, $D$ deferred tax expense (negative for a benefit), and $P$ net cash
taxes paid, all for the same annual period and in millions of dollars. Define
$U$ and $V$ as the residuals needed to connect the reported amounts:
\begin{equation}
 \Delta T=C-P+U,\qquad \Delta N=-D+V.
 \label{eq:apple-tax-bridges}
\end{equation}
These are diagnostic definitions, not assertions that every tax balance has
the same recognition or settlement scope. The first relation gives
$U=11{,}579-6{,}678=4{,}901$. Setting that residual to zero would produce the
wrong closing payable. Assigning it to currency translation, uncertain tax
positions or another journal would require the missing transaction and
classification evidence.

The deferred balance illustrates the second relation. At September 2024,
gross deferred tax assets of 34,873 less a valuation allowance of 8,866 leave
26,007. Deducting deferred tax liabilities of 6,805 gives net deferred tax
assets of 19,202, compared with 17,251 a year earlier
\citep[Note 7, printed p.~41]{apple2024}. The increase is 1,951, so
$V=1{,}951-3{,}031=-1{,}080$. The deferred benefit in profit does not explain
the whole net-balance movement. Its allocation among currency, equity, other
comprehensive income and other causes remains unresolved here.

Note 6's 19,499 deferred-tax-asset presentation is also different from Note 7's
19,202 net position. Their difference of 297 is an arithmetic consequence,
not a separately disclosed liability line or authority to offset jurisdictions.
The reported temporary differences---including research expenditure, deferred
revenue, leases and depreciation---need their own tax bases and recovery
conditions. The note reports indefinitely available Irish foreign-tax and
California R\&D credit carryforwards, but also a valuation allowance for those
credits and some other differences. An indefinite term does not establish that
an asset is recoverable \citep[Notes 6--7, printed pp.~39--41]{apple2024}.

\subsection{A recognised charge, restricted assets and a later payment}

Apple attributes its large 2024 State Aid charge to the September 2024 court
decision confirming the European Commission's recovery decision. This is the
issuer's account of the dispute and its accounting consequence. Note 7 describes
a 15.8 billion dollar payable to Ireland, partly offset by a 4.8 billion US
foreign tax credit and an 823 million decrease in unrecognised tax benefits
\citep[Note 7, printed pp.~39--40]{apple2024}.

The narrative calls the net charge 10.2 billion; the statutory-rate reconciliation
gives the more precise amount 10,246 million. Subtracting the displayed narrative
centres gives $15{,}800-4{,}800-823=10{,}177$, which is 69 million lower.
That is compatible with the coarser inputs: if the two billion-dollar amounts
are rounded to the nearest 0.1 billion and 823 to the nearest million, the
component calculation lies in the closed outer interval
$[10{,}076.5,10{,}277.5]$ million, containing 10,246. The 10.2 billion charge
is compatible with it as well. This is a conditional rounding check, with the
issuer's precise rule and tie convention unverified; the centre difference
does not establish a conflicting source assertion.

The tax charge does not make the settlement funds disappear at recognition.
At 28 September 2024 the escrow comprised 2.6 billion dollars of cash and cash
equivalents and 13.2 billion of current marketable securities, restricted from
general use \citep[Notes 4 and 7, printed pp.~37 and 39]{apple2024}.
The annual liquidity discussion separately identifies the obligation and the
funds held to settle it. The following quarter reports both a payment and the
remaining obligation \citep[printed p.~26]{apple2024}
\citep[printed p.~16]{apple2025q1}.

\begin{center}
\begin{tabularx}{\linewidth}{@{}Xrr@{}}
\toprule
State Aid obligation and settlement & EUR billion & USD billion\\\midrule
Obligation at 28 Sep 2024 & 14.2 & 15.8\\
Payment during the following quarter & 11.0 & 11.9\\
Obligation at 28 Dec 2024 & 3.3 & 3.4\\\bottomrule
\end{tabularx}
\end{center}

The three rows are different dated quantities. In dollars, a movement equation
with opening obligation 15.8, payment 11.9 and closing obligation 3.4 requires
an additional net movement of $3.4+11.9-15.8=-0.5$ billion. Under nearest
rounding to 0.1 billion, its outer interval is $[-650,-350]$ million, excluding
zero. The same calculation in euros gives $3.3+11.0-14.2=0.1$ billion, with
outer interval $[-50,250]$ million, which includes zero. Neither calculation
supplies a common exchange rate or an exact exchange-gain entry. The interim
tax discussion mentions foreign-currency revaluation effects related to the
decision, but does not allocate this 500 million dollar residual
\citep[printed p.~16]{apple2025q1}.

The remaining funds also retain their asset classifications. At December 2024,
restricted cash and cash equivalents were 126 million dollars and restricted
current marketable securities 3.3 billion \citep[Note 4, printed p.~8]{apple2025q1}.
Their displayed sum, 3,426 million, is compatible with the rounded 3.4 billion
obligation. It is not all cash available for ordinary operations. The opening
2.6 billion of restricted cash alone could not have funded the reported
11.9 billion payment; securities sales, maturities, transfers and valuations
need their own histories.

For an isolated settlement illustration, suppose 11,900 million of restricted
cash is available and the amount paid equals the liability's carrying value
at settlement. Debit tax payable and credit restricted cash by 11,900. Assets
and liabilities each fall by 11,900; equity and current-period tax expense do
not change in this posting. This assumes the required cash conversion has
already occurred and omits any associated gain, loss or fee. It explains
settlement of an existing obligation without pretending to reconstruct Apple's
funding and currency entries. The quarter separately reports tax expense of
6,254 and net cash taxes of 18,651 million
\citep[printed pp.~1 and 5]{apple2025q1}; neither should be replaced with the
Ireland payment or treated as another copy of the preceding year's charge.

\subsection{Uncertain tax positions do not supply a payment distribution}

The last part of Note 7 preserves another population: gross unrecognised tax
benefits. Its 2024 movement is $19{,}454+1{,}727-386+2{,}542-1{,}070-229
=22{,}038$ million. The increases concern prior- and current-year positions;
the decreases concern changed prior-year positions, settlements and expired
limitation periods. Interest and penalties are excluded from this rollforward
\citep[Note 7, printed p.~41]{apple2024}. The 1,070 settlement-related decrease
is a movement in that gross balance, not an independently identified total of
cash paid to tax authorities.

Of the rounded 22.0 billion gross amount, the issuer says that 10.8 billion,
if recognised, would affect its effective tax rate. It also describes a
reasonably possible decrease of approximately 5--13 billion over the next
twelve months, primarily concerning transfer pricing and deemed repatriation
tax, while noting uncertainty about examination timing. These are different
conditional statements. The range is neither a probability law nor a promised
cash payment or expense release. A forecast still needs jurisdictional outcomes,
recognition, interest, penalties and settlement dates. Retaining those unknowns
keeps the tax note connected to the company's funding needs without turning
management's dated assessment into an invented future transaction.


\hypertarget{returning-to-the-company-forecast}{%
\section{Returning to the company forecast}\label{method-corporate-methods-returning-to-the-company-forecast}}

The Northstar case remains the broader ordinary-company teaching baseline: collections, purchases/payments, inventory, payroll, PPE, financing, tax, statements, annual/interim and revision views are already specified. Its 2024 collection/purchase/payment solution is 245/155/145. The 2025 forecast in the Northstar reference fixture opens with cash 74.6 and equity 90.6, spends 20 on capital expenditure at period end, pays dividend 8 and uses a base floating rate 7\%. Its simplified tax is settled in the year earned. The running case uses fixed-rate share $\alpha$; here floating-rate share $x=1-\alpha$ describes the same uncertainty. For 2025 the same unknown floating fraction $x\in[0,1]$ determines every
statement amount:
\begin{align*}
 I(x)&=5+2x, & \text{pretax profit}(x)&=82-2x,\\
 \text{tax}(x)&=16.4-0.4x, & \text{profit}(x)&=65.6-1.6x,\\
 \text{operating cash}(x)&=48.6-1.6x, & \text{closing cash}(x)&=95.2-1.6x,\\
 \text{equity}(x)&=148.2-1.6x, & \text{assets}(x)&=288.2-1.6x.
\end{align*}
Receivables 55, inventory 40, net PPE 98, payables 40 and debt 100 complete the
balance sheet. Liabilities are 140 for every $x$, so the common cash/equity
slope preserves the accounting identity. There is no default distribution on
$x$. At a 9\% floating rate the all-floating history has interest 9, profit
62.4, cash 92, equity 145 and assets 285. The annual cash bounds still cannot
establish daily payment solvency without event-date cash flows.

The lease/tax example extends the interaction evidence without changing the original Northstar fixture or pretending that its teaching 20\% tax formula is an IAS 12 method. In a real-company package, each required method is admitted separately and unsupported operations remain visible. The extended obligations for pensions, impairment, provisions, valuation and share-based payment are in the additional corporate methods that follow; omission from a finite company case never proves the actual company has none.
